\scanpage{39}
\mspage{15}

Stratification \deleted{\unsure{équisingulière}}\inserted{équisingulière}
de type \(I\) (\(I\) ens. ordonné fini) : présentation standart.
On suppose donnée une ``catégorie topologique modérée'' \(\mathcal M\)
(\unsure{en} un sens non précisé, pour que ça puisse être aussi celle des
espaces topologiques \(\cdots\)). La présentation est donnée
par un \inserted{contr.} foncteur
\[
 \begin{aligned}
 \Sigma:\ \operatorname{Drap}(I)&\longrightarrow\mathcal M\\
 (i_0,\ldots,i_n)&\longmapsto\Sigma_{i_0\ldots i_n}
 \end{aligned}
\]
avec les propriétés suivantes

(1) Le \deleted{\unsure{app}} morphisme
\[
 \Sigma_{i_0\ldots i_n}\longrightarrow\Sigma_{i_0,\ldots,i_p}
\]
(relatif à une inclusion de drapeaux
\(d\unsure{\hookrightarrow}d'\) telle
que \(d\) soit un segment initial de \(d'\)) est \uline{fibrant}\(^*\).

\margin{NB \inserted{on} \unsure{ne poser}
cette condition que dans le cas équisingulier.}

\margin{\(^*\) Cela est trop faible, $=$ pour
l'équisingularité --- il faut
supposer que pour \(d\) fixé, et pour
\(d'\unsure{\supset|=}d\amalg\delta\) \unsure{variable}
(où \(\delta\) est \unsure{donc}
un drapeau quelconque de
l'ens \(I'\) des \(i\in I\) \unsure{fort\(^{t}\)|fortement} \uline{\unsure{supérieur|supérieurs}}
\unsure{à} \(\unsure{\max} i_p\))
la \deleted{\unsure{famille}}\unsure{stratification}
qu'ils forment au dessus de
\(\Sigma_{d'}\) est loc. triviale.}

\deleted{(2)} (\deleted{\unsure{Enfin}} il suffit de le supposer pour
\[
 \Sigma_{i_0\ldots i_n}\longrightarrow\Sigma_{i_0\ldots i_{n-1}}
\]
)

(2) Le morphisme
\[
 \deleted{\unsure{\Sigma_{d'}}\longrightarrow\Sigma_d}
 \qquad
 \Sigma_{i_0\ldots i_n}
 \longrightarrow
 \Sigma_{i_0\ldots\widehat{i_p}\ldots i_n},
 \qquad \boxed{p\ne n}
\]
est un plongement-bord
\inserted{i.e. plongement \inserted{\unsure{bord}} cylindrique}
(NB si \(p=n\), c'est un morphisme fibrant.)

\margin{On va supposer
de plus \unsure{données|donnés}
une \deleted{\unsure{famille}} \inserted{des}
structures cylindriques
correspondantes, et \unsure{formuler}
des conditions plus fortes que
(2), \unsure{énoncées en langage}
\unsure{suivant} ---}

\uline{Cor.} Soit \(d\unsure{\hookrightarrow}d'\)
tel que \(d\) et \(d'\) aient \(\unsure{=}\)
\unsure{objt maximals}, alors
\(\Sigma_{d'}\longrightarrow\Sigma_d\) est un plongement.

\deleted{\unsure{(3)}}

\margin{38}
